\subsection{Είδη LCS}

Οι παρακάτω κατηγορίες και περιγραφές έχουν παρθεί από το συμπληρωματικό υλικό της έρευνας των \cite{Garcia2022}:
\subsubsection{Ηλεκτροχημικοί αισθητήρες (Electrochemical sensors ,EC)}

Οι ηλεκτροχημικοί αισθητήρες μετατρέπουν την επίδραση μιας ηλεκτροχημικής αλληλεπίδρασης μεταξύ του επιλεγμένου για μέτρηση  αερίου και του ηλεκτροδίου σε ηλεκτρικό σήμα, το οποίο είναι ανάλογο της συγκέντρωσης του αερίου [1, 2]. Οι αισθητήρες αυτοί περιλαμβάνουν ένα στοιχείο που περιέχει ηλεκτρολύτη καθώς και δύο ή τρία ηλεκτρόδια, και όταν ο αισθητήρας έρχεται σε επαφή με το επιλεγμένο αέριο, πραγματοποιείται οξείδωση ή μείωση, παράγοντας ένα μετρήσιμο ρεύμα μεταξύ των ηλεκτροδίων εργασίας και αντίθεσης. Οι περισσότεροι LCS διαθέτουν ένα τρίτο αναφορικ» ηλεκτρόδιο προκειμένου να βελτιώσουν τη γραμμικότητα και την ευαισθησία [2, 3]. Οι αισθητήρες αυτοί μπορεί να είναι ευάλωτοι σε παρεμβολές από περιβαλλοντικές παραμέτρους και από άλλες χημικές ενώσεις [4].


\subsubsection{Ημιαγωγοί Μεταλλικού Οξειδίου (Metal Oxide Semiconductors,MOS)}

Η αρχή λειτουργίας των αισθητήρων ημιαγωγικού μεταλλικού οξειδίου 
βασίζεται στην απορρόφηση οξυγόνου σε υψηλές θερμοκρασίες στην επιφάνεια των οξειδίων μετάλλων,
 παγιδεύοντας έτσι ηλεκτρόνια από το εσωτερικού του υλικού, αυξάνοντας ή μειώνοντας την αντίσταση του αισθητήρα για υλικά τύπου n ή p,
  αντίστοιχα. Στον αέρα, σε θερμοκρασίες μεταξύ $150^\circ \mathrm{C}$ και $400^\circ \mathrm{C}$, το οξυγόνο απορροφάται στην επιφάνεια.
   Το επιλεγμένο προς μέτρηση αέριο  αντιδρά με το προσροφημένο οξυγόνο ή άμεσα με το οξείδιο, προκαλώντας μεταβολή της αντίστασης του αισθητήρα, η οποία ανιχνεύεται ως σήμα του αισθητήρα και σχετίζεται με τη συγκέντρωση του  επιλεγμένο προς μέτρηση αερίου [5, 6].
    Οι αισθητήρες MOS αποτελούν τη χαμηλή κόστους λύση για την παρακολούθηση της ποιότητας του αέρα, ωστόσο παρουσιάζουν σημαντικές προκλήσεις στην ποσοτικοποίηση, όπως διασταυρούμενες ευαισθησίες (cross-sensitivity)  και εξάρτηση από την σχετική υγρασία και τη θερμοκρασία [7].

\subsubsection{Μη Διασπειρόμενο Υπέρυθρος,(Non-Dispersive InfraRed ,NDIR)}

Αυτός ο τύπος αισθητήρα βασίζεται στις οπτικές ιδιότητες του επιλεγμένου προς μέτρηση αερίου, δηλαδή στα μήκη κύματος στα οποία το αέριο απορροφά το φως. Ο αισθητήρας αποτελείται από μια πολυχρωματική πηγή φωτός, συνήθως μια λυχνία ή LED, ένα κυψελίδιο αερίου, ένα φίλτρο ζώνης (band-pass) για τον περιορισμό του εισερχόμενου φωτός στη συγκεκριμένη φασματική ζώνη απορρόφησης του αναλυτή, και τον ανιχνευτή. Πολλά αέρια ρύπων είναι ενεργά στο υπέρυθρο (IR) φάσμα (π.χ. CO, CO\textsubscript{2}, SO\textsubscript{2}, NO\textsubscript{x}, N\textsubscript{2}O, N\textsubscript{3}, HCl, HF και CH\textsubscript{4}), και μπορούν δυνητικά να παρακολουθηθούν με αυτή τη μέθοδο. Η ποσοτικοποίηση στις μετρήσεις NDIR βασίζεται στον νόμο απορρόφησης Beer-Lambert [9, 10].

\subsubsection{Ανιχνευτής Φωτοϊονισμού (Photoionization detectors, ID)}

Η λειτουργία τους βασίζεται στην απορρόφηση φωτονίων από μια πηγή φωτοεκκένωσης από ένα μόριο του στοχευόμενου αερίου, οδηγώντας στο σχηματισμό ενός διεγερμένου μορίου, δηλαδή φωτοϊόντων. Όταν η ενέργεια των φωτονίων hν υπερβαίνει την ενέργεια ιονισμού του μορίου, εκπέμπεται ένα ηλεκτρόνιο και σχηματίζεται ένα κατιόν, το οποίο κατευθύνεται προς ένα ηλεκτρόδιο, παράγοντας ηλεκτρικό ρεύμα που εξαρτάται από τον αριθμό των διεγερμένων μορίων και, συνεπώς, από τη συγκέντρωση του ανιχνευόμενου αερίου [8]. Στους PID χρησιμοποιούνται διάφοροι τύποι λυχνιών. Οι αισθητήρες PID εμφανίζουν σύντομο χρόνο απόκρισης (λίγα δευτερόλεπτα) και δεν είναι ειδικοί για κάποιο συγκεκριμένο αναλυτή, αλλά παρουσιάζουν υψηλή ευαισθησία σε αλκοόλες, αρωματικές ενώσεις, αλδεΰδες και κετόνες, λιπαρά οξέα και πολυκυκλικούς αρωματικούς υδρογονάνθρακες (PAH), με τους τελευταίους να ανιχνεύονται ακόμη και σε χαμηλότερες συγκεντρώσεις σε σχέση με αυτές που μπορούν να επιτευχθούν με αισθητήρες NDIR [9].



\subsubsection{Διασκορπισμός φωτός (Light Scattering)}

Στην περίπτωση των σωματιδίων αιωρούμενης ύλης PM, η βασική αρχή ανίχνευσης στηρίζεται στη σκέδαση, ή διασκορπισμό φωτός, κατά την οποία ο αέρας από το περιβάλλον που παρακολουθείται οδηγείται μέσα στο περίβλημα του αισθητήρα, όπου χρησιμοποιείται μια πηγή φωτός σε συνδυασμό με φωτοδίοδο. Τα σωματίδια προκαλούν σήμα σκέδασης φωτός, το οποίο συνήθως μετατρέπεται σε αναλογικό σήμα τάσης ή σήμα διαμόρφωσης πλάτους παλμών (pulse width modulation), που χρησιμοποιείται για την αναπαράσταση των μετρούμενων συνθηκών, όπως πλήθος παλμών, ύψος ή διάρκεια παλμών. Το σήμα αυτό μπορεί να αναλυθεί ως προς τη συχνότητα των παλμών, εκτιμώντας έτσι τον αριθμό των σωματιδίων, και ως προς το αντίστοιχο ύψος των παλμών, παρέχοντας εκτίμηση του μεγέθους των σωματιδίων. Εναλλακτικά, η συνολική ένταση του σκορπισμένου φωτός μπορεί να χρησιμοποιηθεί για την εκτίμηση ενός συνόλου σωματιδίων [11]. Η μέθοδος αυτή ονομάζεται νεφελομετρική σκέδαση.

\subsubsection{Θάλαμοι ιονισμού (Ionization Chambers)}

θαλάμη ιονισμού αποτελεί ενεργό στοιχείο αισθητήρα που χρησιμοποιείται συχνά για την ανίχνευση ιοντίζουσας ακτινοβολίας. Στην περίπτωση του ραδονίου, χρησιμοποιούνται για την ανίχνευση σωματιδίων α και μπορούν να λειτουργούν για αξιολόγηση βραχυχρόνιας ή μακροχρόνιας έκθεσης. Όταν ένα φορτισμένο σωματίδιο διέρχεται από τη θαλάμη παρουσία ιοντίζουσας ακτινοβολίας, παράγονται ζεύγη ηλεκτρονίου-ιόντος, τα οποία επιταχύνονται λόγω του εφαρμοζόμενου ηλεκτρικού πεδίου μέχρι να συγκρουστούν με τα ηλεκτρόδια [12]. Στη συνέχεια, ένα ηλεκτρονικό κύκλωμα παράγει έναν ηλεκτρικό παλμό ανά σύγκρουση, οι οποίοι στη συνέχεια καταμετρώνται για συγκεκριμένα χρονικά διαστήματα [13]. Συνήθως, οι LCS για το ραδόνιο χρησιμοποιούν θαλάμες ιονισμού [14, 15], λόγω της απλής κατασκευής τους, της προσιτής τιμής και της ικανοποιητικής ακρίβειας.

\subsubsection{Φασματομετρία άλφα (Alpha Spectrometry )}

Η άλφα-φασματομετρία μπορεί να χρησιμοποιηθεί για την ταυτοποίηση και ποσοτικοποίηση των σωματιδίων α που εκπέμπονται κατά τη διαδικασία διάσπασης ραδιονουκλεϊδίων, όπως το ραδόνιο. Στη διαδικασία αυτή, η ιοντίζουσα ακτινοβολία μπορεί να ανιχνευθεί μέσω αλφα-φασματομετρίας [16, 17]. Συνήθως, οι υλοποιήσεις χαμηλού κόστους της αλφα-φασματομετρίας βασίζονται σε φωτοδιόδους PIN, ευαίσθητες σε σωματίδια α [18], οι οποίες μπορούν να υλοποιηθούν χρησιμοποιώντας εμπορικά διαθέσιμα εξαρτήματα [19].
